\documentclass[10pt,a4paper]{amsart}
\usepackage{amsmath,amssymb,amsthm,graphicx}
% Retained minimal preamble for the Stacks Project chapter more-morphisms.tex (More on Morphisms).
% Verbatim excerpt of the chapter preamble: Xy-pic setup (native preamble lines 14-18) and the native
% theorem declarations with the mathematical macros (native preamble lines 154-217). Copied mechanically;
% nothing added, removed or redefined. The wrapper class and the amsmath/amssymb/amsthm/graphicx base
% packages are supplied by the assembler, so the class/externaldocument part of the native preamble is omitted.
% For commutative diagrams we use Xy-pic
\usepackage[all]{xy}

% We use 2cell for 2-commutative diagrams.
\xyoption{2cell}
% Theorem environments.
%
\theoremstyle{plain}
\newtheorem{theorem}[subsection]{Theorem}
\newtheorem{proposition}[subsection]{Proposition}
\newtheorem{lemma}[subsection]{Lemma}

\theoremstyle{definition}
\newtheorem{definition}[subsection]{Definition}
\newtheorem{example}[subsection]{Example}
\newtheorem{exercise}[subsection]{Exercise}
\newtheorem{situation}[subsection]{Situation}

\theoremstyle{remark}
\newtheorem{remark}[subsection]{Remark}
\newtheorem{remarks}[subsection]{Remarks}

\numberwithin{equation}{subsection}

% Macros
%
\def\lim{\mathop{\mathrm{lim}}\nolimits}
\def\colim{\mathop{\mathrm{colim}}\nolimits}
\def\Spec{\mathop{\mathrm{Spec}}}
\def\Hom{\mathop{\mathrm{Hom}}\nolimits}
\def\Ext{\mathop{\mathrm{Ext}}\nolimits}
\def\SheafHom{\mathop{\mathcal{H}\!\mathit{om}}\nolimits}
\def\SheafExt{\mathop{\mathcal{E}\!\mathit{xt}}\nolimits}
\def\Sch{\mathit{Sch}}
\def\Mor{\mathop{\mathrm{Mor}}\nolimits}
\def\Ob{\mathop{\mathrm{Ob}}\nolimits}
\def\Sh{\mathop{\mathit{Sh}}\nolimits}
\def\NL{\mathop{N\!L}\nolimits}
\def\CH{\mathop{\mathrm{CH}}\nolimits}
\def\proetale{{pro\text{-}\acute{e}tale}}
\def\etale{{\acute{e}tale}}
\def\QCoh{\mathit{QCoh}}
\def\Ker{\mathop{\mathrm{Ker}}}
\def\Im{\mathop{\mathrm{Im}}}
\def\Coker{\mathop{\mathrm{Coker}}}
\def\Coim{\mathop{\mathrm{Coim}}}

% Boxtimes
%
\DeclareMathSymbol{\boxtimes}{\mathbin}{AMSa}{"02}

%
% Macros for moduli stacks/spaces
%
\def\QCohstack{\mathcal{QC}\!\mathit{oh}}
\def\Cohstack{\mathcal{C}\!\mathit{oh}}
\def\Spacesstack{\mathcal{S}\!\mathit{paces}}
\def\Quotfunctor{\mathrm{Quot}}
\def\Hilbfunctor{\mathrm{Hilb}}
\def\Curvesstack{\mathcal{C}\!\mathit{urves}}
\def\Polarizedstack{\mathcal{P}\!\mathit{olarized}}
\def\Complexesstack{\mathcal{C}\!\mathit{omplexes}}
% \Pic is the operator that assigns to X its picard group, usage \Pic(X)
% \Picardstack_{X/B} denotes the Picard stack of X over B
% \Picardfunctor_{X/B} denotes the Picard functor of X over B
\def\Pic{\mathop{\mathrm{Pic}}\nolimits}
\def\Picardstack{\mathcal{P}\!\mathit{ic}}
\def\Picardfunctor{\mathrm{Pic}}
\def\Deformationcategory{\mathcal{D}\!\mathit{ef}}
\usepackage{cleveref}
\crefname{originallemma}{original lemma}{original lemmas}
\Crefname{originallemma}{Original lemma}{Original lemmas}
\makeatletter
\newlabel{lemma-etale-makes-quasi-finite-finite-at-point}{{02LK}{source}{}{sourcepointer-1}{}}
\newlabel{lemma-etale-makes-quasi-finite-finite-at-point@cref}{{[originallemma][0][]02LK}{[1][0][]source}{}{}{}}
\makeatother
\crefname{originallemma}{original lemma}{original lemmas}
\Crefname{originallemma}{Original lemma}{Original lemmas}
\makeatletter
\newlabel{lemma-etale-splits-off-quasi-finite-part-technical-variant}{{02LO}{source}{}{sourcepointer-2}{}}
\newlabel{lemma-etale-splits-off-quasi-finite-part-technical-variant@cref}{{[originallemma][0][]02LO}{[1][0][]source}{}{}{}}
\makeatother
\crefname{originallemma}{original lemma}{original lemmas}
\Crefname{originallemma}{Original lemma}{Original lemmas}
\makeatletter
\newlabel{lemma-weighting-check-after-etale-base-change}{{0F39}{source}{}{sourcepointer-3}{}}
\newlabel{lemma-weighting-check-after-etale-base-change@cref}{{[originallemma][0][]0F39}{[1][0][]source}{}{}{}}
\makeatother
\title{Composition of weightings for locally quasi-finite morphisms}
\author{The Stacks Project Authors}
\date{stacks-project:more-morphisms.tex:local-snapshot-20261011; GNU FDL 1.2 or later}
\begin{document}
\maketitle
\noindent\textit{Editorial note.} Counted span is the author's own proof environment, source lines 22261-22309, immediately following the statement at 22249-22259.\par\medskip
\noindent\textit{Editorial note.} The proof cites three chapter lemmas by \textbackslash{}ref; these are navigation pointers to standing inputs of the surrounding chapter, so their statements are not re-counted here.\par\medskip
\noindent\textit{Editorial note.} The closing line 'It follows immediately that ...' is retained verbatim: it only multiplies the two constant values already produced inside the span.\par\medskip
\noindent\textit{Allowed standard prerequisites (editorial).}\begin{itemize}
\item Criterion for a weighting after etale base change. chapter lemma lemma-weighting-check-after-etale-base-change, cited at proof lines 22273-22276. Lets the local constancy of the integral be checked on an etale neighbourhood of a point of the base.
\item Etale splitting off of the quasi-finite part. chapter lemma lemma-etale-splits-off-quasi-finite-part-technical-variant, cited at proof lines 22280-22281. Decomposes the finite part over an etale neighbourhood into pieces each having a unique point with purely inseparable residue extension over the base point.
\item Quasi-finiteness becomes finiteness after an elementary etale neighbourhood. chapter lemma lemma-etale-makes-quasi-finite-finite-at-point, cited at proof line 22293. Produces an open neighbourhood of the image of the marked point over which the base change is finite, yielding the intermediate finite morphism V -> W.
\item Etale morphisms are open. standard Stacks fact used at proof line 22274. Justifies replacing the base by an etale neighbourhood when testing local constancy.
\end{itemize}\medskip
\ifdefined\hypertarget\hypertarget{sourcepointer-1}{}\fi
\noindent\textit{Source reference.} Original lemma 02LK: sources/stacks-more-morphisms/pinned-original-source.tex, source lines 11924--11924. Navigation pointer only: the counted proof cites this statement by \textbackslash{}ref as a standing input; the deduction the counted proof has to make is carried out inside the counted span, and no part of that deduction is delegated to this pointer.\par\medskip
\ifdefined\hypertarget\hypertarget{sourcepointer-2}{}\fi
\noindent\textit{Source reference.} Original lemma 02LO: sources/stacks-more-morphisms/pinned-original-source.tex, source lines 12136--12136. Navigation pointer only: the counted proof cites this statement by \textbackslash{}ref as a standing input; the deduction the counted proof has to make is carried out inside the counted span, and no part of that deduction is delegated to this pointer.\par\medskip
\ifdefined\hypertarget\hypertarget{sourcepointer-3}{}\fi
\noindent\textit{Source reference.} Original lemma 0F39: sources/stacks-more-morphisms/pinned-original-source.tex, source lines 22111--22111. Navigation pointer only: the counted proof cites this statement by \textbackslash{}ref as a standing input; the deduction the counted proof has to make is carried out inside the counted span, and no part of that deduction is delegated to this pointer.\par\medskip
\begin{lemma}
\label{lemma-weighting-composition}
Let $f : X \to Y$ and $g : Y \to Z$ be locally quasi-finite morphisms.
Let $w_f : X \to \mathbf{Z}$ be a weighting of $f$ and let
$w_g : Y \to \mathbf{Z}$ be a weighting of $g$. Then the function
$$
X \longrightarrow \mathbf{Z},\quad
x \longmapsto w_f(x) w_g(f(x))
$$
is a weighting of $g \circ f$.
\end{lemma}

\begin{proof}
Let us set $w_{g \circ f}(x) = w_f(x) w_g(f(x))$ for $x \in X$.
Consider a diagram
$$
\xymatrix{
X \ar[d]_{g \circ f} & U \ar[l] \ar[d]^\pi \\
Z & W \ar[l]
}
$$
where $W \to Z$ is \'etale, $U \subset X_W$ is open, and $U \to W$ finite.
We have to show that $\int_\pi w_{g \circ f}|_U$ is locally constant.
Choose a point $w \in W$. By
Lemma \ref{lemma-weighting-check-after-etale-base-change}
(and the fact that \'etale morphisms are open) it suffices to show that
$\int_\pi w_{g \circ f}|_U$ is constant after replacing
$(W, w)$ by an \'etale neighbourhood.
After replacing $(W, w)$ by an \'etale neighbourhood we may
assume $U = U_1 \amalg \ldots \amalg U_n$ where each $U_i$
has a unique point $u_i$ lying over $w$ such that
$\kappa(u_i)/\kappa(w)$ is purely inseparable, see
Lemma \ref{lemma-etale-splits-off-quasi-finite-part-technical-variant}.
Clearly, it suffices to show that $\int_{U_i \to W} w_{g \circ f}|_{U_i}$
is constant in an \'etale neighbourhood of $w$.
This reduces us to the case discussed in the next paragraph.

\medskip\noindent
We have $w \in W$ and there is a unique point $u \in U$ lying
over $w$ with $\kappa(u)/\kappa(w)$ purely inseparable.
Consider the point $v = f(u) \in Y$. After replacing
$(W, w)$ by an elementary \'etale neighbourhood we may
assume there is an open neighbourhood $V \subset Y_W$
of $v$ such that $V \to W$ is finite, see
Lemma \ref{lemma-etale-makes-quasi-finite-finite-at-point}.
Then $f_W^{-1}(V) \cap U$ is an open neighbourhood of $u$
where $f_W : X_W \to Y_W$ is the base change of $f$ to $W$.
Hence after Zariski shrinking $W$, we may assume $f_W(U) \subset V$.
Thus we obtain morphisms
$$
U \xrightarrow{a} V \xrightarrow{b} W
$$
and $U \to V$ is finite as $V \to W$ is separated (because finite).
Since $w_f$ and $w_g$ are weightings of $f$ and $g$
we see that $\int_a w_f|_U$ is locally constant on $V$ and
$\int_b w_g|_V$ is locally constant on $W$. Thus after shrinking
$W$ one more time we may assume these functions are constant
say with values $n$ and $m$. It follows immediately that
$\int_\pi w_{g \circ f}|_U = \int_{b \circ a} w_{g \circ f}|_U$
is constant with value $nm$ as desired.
\end{proof}

\end{document}
